\begin{lemma}\label{lem:decic242-case26-cubic}
In case $26$, the degree-$4$ cusp complement contains one closed
cubic cusp and one rational cusp.
\end{lemma}
\begin{proof}
Use the first map to $X_0(42)$ in the exact quotient sieve.
Its normalized noncuspidal class is
\[
 v=(1,3,6,12,0)\in
 G=[2,4,48,48,0].
\]
The rational-image subgroup $H=[2,48]$ has generators whose
first coordinates are all even. The first coordinate modulo $2$
therefore vanishes on $H$, whereas it takes the value $1$ on $v$.
The residue cusp complement cannot be the reduction of a rational
cuspidal divisor: its pushforward would contribute a class in $H$
and could not correct $v$.

The global source cusp signature is $1^{18}2^9 6^6$, and the
signature at $5$ is $1^{18}2^9 3^6 6^3$.
Every rational or quadratic residue cusp is globally intact.
Hence every degree-$4$ cusp complement without a cubic cusp lifts
to a rational cuspidal divisor. This is impossible by the preceding
class test. The only remaining pattern is one cubic cusp and one
rational cusp.
\end{proof}

The exact subgroup vectors and noncuspidal class are in
\path{geometry/x0_42_proven_picard_v2.m} and
\path{geometry/x0_42_full_arithmetic.m}.
The global cusp pushforwards are certified in
\path{geometry/x0_42_direct_cusp_pushforward.m}.
An independent integer enumeration verifies both source cusp
signatures and the intact rational and quadratic cusps; see
\path{case26_cubic_complement.json}.
